\subsection{Overview}
\TD\ identifies candidate protein-coding regions (CDS) within transcript sequences in two stages
(Fig.~\todo{workflow figure}). \emph{LongOrfs} enumerates open reading frames (ORFs) above a minimum
length. \emph{Predict} trains a coding-sequence model on the longest of these ORFs and uses it,
together with a GC-calibrated length criterion and optional homology evidence, to select the likely
coding ORF(s) in each transcript. It then refines start codons and reports CDS, UTRs and proteins.
An optional third stage projects transcript-level CDS onto a reference genome. Unless stated
otherwise, the algorithm is identical in the Perl implementation (\TD\ v6.1) and in \PyTD\
\todo{version}; Section~\ref{sec:impl} describes their concordance.

\subsection{ORF enumeration}
Each transcript is scanned on the forward strand and, unless the data are declared strand-specific
(\texttt{-S}), on the reverse complement. Stop codons come from the selected genetic code. Twenty-four
named codes are supported, including the ciliate (\emph{Tetrahymena}), \emph{Candida}, Euplotid and
twelve mitochondrial codes. Start codons are ATG in all codes.

In each frame, every stop codon is paired with the most upstream in-frame ATG that follows the
preceding in-frame stop, giving one maximal ORF per stop. ORFs may run off either end of the
transcript, which accommodates incomplete assemblies:
\begin{itemize}
  \item An ORF with no in-frame stop upstream of its first codon is \emph{5$'$ partial}. It begins at
        the first full codon of the transcript.
  \item An ORF that reaches the 3$'$ end without a stop is \emph{3$'$ partial}.
\end{itemize}
ORFs are labelled complete, 5$'$ partial, 3$'$ partial or internal (both). ORFs of at least $m$
codons are retained, where $m=100$ by default; the stop codon, or its absence at a 3$'$ end, counts
as one codon. \texttt{--complete\_orfs\_only} restricts the output to ORFs with both a start and a
stop.

\subsection{Background composition}
Nucleotide frequencies $q(b)$, $b\in\{A,C,G,T\}$, are pooled over all transcripts. Without
\texttt{-S}, counts from both strands are combined, so that $q(A)=q(T)$ and $q(C)=q(G)$. The G+C
content used below is $2\,q(C)$.

\subsection{Training set}
The coding model is trained without supervision on long ORFs, which are very unlikely to arise by
chance:
\begin{enumerate}
  \item Candidates are the $10T$ longest ORFs (by default $T=500$). ORFs with CDS longer than 5{,}000\,nt
        are excluded.
  \item Near-identical proteins are removed greedily, longest first:
        \begin{itemize}
          \item Each candidate protein is cut into consecutive, non-overlapping 5-residue words.
          \item If its distinct words make up less than 30\% of all its words, it is skipped as low
                complexity.
          \item It is discarded if more than 80\% of its distinct words occur in a single
                previously retained protein. Words of retained proteins are indexed at every offset.
        \end{itemize}
  \item The $T$ longest surviving ORFs form the training set. They may be partial and from either
        strand.
\end{enumerate}

\subsection{Coding potential}
Coding sequence is modelled as a three-periodic (inhomogeneous) fifth-order Markov chain
\citep{Borodovsky1993,Parra2000}. Let $N_f(w)$ be the number of training-set occurrences of the
$k$-mer $w$ ($k\le 6$) whose last base lies at codon position $f\in\{0,1,2\}$. The transition
probability of base $x_i$ in frame $f=i \bmod 3$ given the preceding $\min(i,5)$ bases is
\begin{equation}
  P(x_i \mid x_{i-k..i-1}, f) = \frac{N_f(x_{i-k..i}) + 1}{N_{(f-1)\bmod 3}(x_{i-k..i-1}) + 4}.
\end{equation}
Each position contributes the log-odds $\ell_i=\ln\!\left[P(x_i\mid\cdot,f)/q(x_i)\right]$ against the
zeroth-order background. An ORF sequence $s$ is scored as $S(s)=\sum_i \ell_i$, excluding a
terminal stop codon of the selected genetic code.

Six scores are computed for every candidate:
\begin{itemize}
  \item $S_1$ for the ORF in its own frame;
  \item $S_2$, $S_3$ for the same sequence offset by one and two nucleotides;
  \item $S_4$--$S_6$ for the three frames of its reverse complement.
\end{itemize}

\subsection{ORF selection}
An ORF is a candidate coding region if any of the following hold:
\begin{enumerate}
  \item it is supported by homology evidence;
  \item its CDS length is at least $L^\ast(\mathrm{GC})$ (see below);
  \item $S_1>0$ and $S_1>\max(S_2,S_3)$.
\end{enumerate}
Rule~3 compares the reading frame with the two alternative forward frames. An earlier rule also
required $S_1$ to exceed the reverse-complement scores. It was dropped in v5 because coding
sequence on the opposite strand of a gene often scores highly. \todo{cite the 2018 decision-tree
analysis, Supplementary Fig.}

\paragraph{GC-calibrated long-ORF rule.} Long ORFs are kept even when they score poorly. This
recovers genes with atypical composition, such as horizontally transferred genes or genes in
transcriptomes containing several organisms. $L^\ast$ is the 0.999 quantile of the lengths of ORFs
that LongOrfs extracts from random sequence of the same G+C content. The random sequences are
i.i.d., scanned on both strands, with partial ORFs allowed and a 300-nt minimum. $L^\ast$ was
estimated by simulation in 5\% GC steps between 25\% and 80\%, with at least 85{,}000 random ORFs per
level (Supplementary Table~\todo{S1}). It rises from 519\,nt at 25\% GC to 2{,}436\,nt at 80\%. The
rule is disabled above 80\% GC. A fixed length can be set instead
(\texttt{--retain\_long\_orfs\_mode strict}).

\paragraph{Homology.} Optionally, the proteins of all candidate ORFs can be searched against a
protein database with DIAMOND \citep{Buchfink2021} or BLASTP, and against Pfam \citep{Mistry2021}
with HMMER \citep{Eddy2011}. ORFs with hits are retained and ranked first.

\paragraph{Ranking and overlap resolution.} Within each transcript, candidates are ranked by:
\begin{enumerate}
  \item number of homology hits;
  \item $S_1$;
  \item length.
\end{enumerate}
They are then accepted greedily. A candidate is rejected if it overlaps an accepted ORF on either
strand by more than 10\% of either ORF's length; homology-supported ORFs are always accepted.
Several non-overlapping ORFs per transcript can therefore be reported, as in polycistronic or
chimeric assemblies. \texttt{--single\_best\_only} keeps a single ORF per transcript.

\subsection{Start-codon refinement}
\label{sec:start}
A 5$'$-partial ORF begins at the first in-frame codon of the transcript rather than at an ATG. If
the transcript is complete, the true start lies downstream. \TD\ learns the sequence context of
translation start sites from the training set:
\begin{itemize}
  \item Each feature spans 20\,nt upstream of an ATG, the ATG, and 10\,nt downstream.
  \item Positive features are the start codons of training ORFs.
  \item Negative features are all ATGs downstream of those starts.
\end{itemize}
Position weight matrices (pseudocount 0.1) are estimated for positives and negatives. A candidate is
scored as the log-likelihood ratio over a sub-window of $u$ upstream and $d$ downstream positions,
divided by $u+d$.

The procedure has three steps:
\begin{enumerate}
  \item Noisy positives are depleted iteratively, starting from a random 30\% subset.
  \item The sub-window $(u,d)$, with $1\le u\le20$ and $1\le d\le10$, is chosen by maximising the
        cross-validated area under the ROC curve. The decision threshold is chosen by maximising
        the F1 score.
  \item For each ORF that does not begin with ATG, the in-frame ATGs within
        $\max(30\,\mathrm{nt}, 15\%\ \text{of ORF length})$ of its 5$'$ end are scored. The ORF is
        moved to the highest-scoring ATG at or above the threshold, and relabelled complete (or
        3$'$ partial).
\end{enumerate}
\todo{Update with the stabilised procedure if adopted (analyses/start\_refinement\_stability). Note the
2026 fix: earlier releases learned the threshold on length-normalised scores but applied it to
unnormalised scores on a window shifted by one nucleotide.}

\subsection{Outputs and genome projection}
\TD\ writes GFF3 (gene, mRNA, CDS, UTR), BED, and CDS and peptide FASTA files. The record headers
carry ORF type, score and homology hits.

For transcripts reconstructed from a genome alignment (StringTie, Cufflinks, PASA, GMAP, or
long-read isoform pipelines), \texttt{cdna\_alignment\_orf\_to\_genome\_orf.pl} projects the
transcript CDS onto the genome through the exon structure, giving genome-coordinate gene models in
GFF3. ORFs from different isoforms of a gene keep a common gene identifier. ORFs predicted on the
opposite strand of a multi-exon transcript are discarded. For a single-exon transcript, such an ORF
reverses the transcript's orientation.

\subsection{Implementations}
\label{sec:impl}
\TD\ is implemented in Perl (\url{https://github.com/TransDecoder/TransDecoder}). Version~6 adds a
single wrapper that runs genome-based cDNA extraction, LongOrfs, optional DIAMOND/BLASTP and
HMMER searches, Predict, and genome projection.

\PyTD\ is a pure-Python reimplementation with Biopython as its only dependency. It reproduces the
Perl algorithm, including the ordering of ties. On five test data sets the two implementations
produce identical ORFs, coding scores and selections; they differ only in start codons chosen by the
randomised PWM training (Section~\ref{sec:start}; Supplementary Table~\todo{S2}). \PyTD\ runs
\todo{2.7}$\times$ faster \todo{(benchmark runtime table)}.
